\scanpage{276}
% archive foliation: 275

\begin{center}
Rapport de la Commission relative au

recrutement des enseignants
\end{center}

I. Organisation du Département de Mathématique.

Le \uline{Département de Mathématique} de la Faculté des Sciences d'Orsay se compose d'une \uline{Ecole de Mathématique} et d'un \uline{Laboratoire de Mathématique}

1°/ L'\uline{Ecole de Mathématique} a la responsabilité des enseignements de premier et second cycle, y compris les enseignements non classiques qu'elle a la charge de développer en collaboration avec les autres départements de la Faculté.

2°/ Le \uline{Laboratoire de Recherche Mathématique} a la responsabilité des enseignements de troisième cycle et, d'une façon générale, de toutes les activités ayant pour but de faire de l'ensemble du Département un milieu favorable à la production mathématique.

Tout enseignant du Département de Mathématique, à un moment déterminé, appartient à l'une ou l'autre des deux sections. Les qualifications professionnelles requises pour être membre permanent de l'une ou l'autre des deux sections sont différentes.

3°/ Les décisions concernant l'organisation de l'enseignement, les nominations à des postes d'enseignement ou de recherche, les éventuels changements organiques à apporter au département de Mathématique, sont prises par l'\uline{Assemblée du département de Math.} de la Fac. des Sc.
